% 【本文件写什么】impl 实现层：dummy
% - 定位：平台驱动占位。
% - crate 路径：os/components/wateros-driver/driver-impl/impl-dummy/src/
% - 如何实现对应 api-v0 trait：关键类型、状态机、数据结构
% - 算法或硬件交互流程，可用伪代码/流程图
% - 本 impl 启用的 Cargo [features] 与 arch feature，impl-riscv64 等
% - 与其它 impl 变体的差异、切换 feature 时的影响
% - 性能/内存假设与已知限制
% 【事实来源】
% - 上述 crate 源码与 Cargo.toml
% - os/feature-tree.txt 中指向本 impl 的 feature 链

\subsubsection{impl-dummy，platform}

\code{wateros-driver-impl-dummy}为无硬件目标的构建占位：不解析DTB、不实现\code{init\_when\_boot}/\code{init\_after\_boot}/\code{physical\_ram\_end\_exclusive}等平台入口；聚合层对 \code{impl-dummy} 的 \code{init\_when\_boot} 为空操作，\code{init\_after\_boot} 不调用平台路径。

导出\code{add(left, right)}占位函数维持依赖图；\code{driver::test()}在\code{impl-dummy} feature下跳过QEMU探测并打\code{info}日志。非QEMU板级若需真实bring-up，应新增独立\code{impl-*} crate并在根\code{Cargo.toml} feature中接线，而非扩展dummy。
